\documentclass[11pt]{article}
\usepackage[a4paper,margin=1in]{geometry}
\usepackage{amsmath,amssymb,booktabs,longtable,array,multirow}
\usepackage{caption}
\usepackage{enumitem}

\begin{document}

\title{Problem Set 4: Derivatives}
\author{}
\date{}
\maketitle

\section*{Exercise 2. Stochastic Calculus and Pricing}
\subsection*{2.1 From random walk to Brownian motion.}

Let $\{\omega_i\}_{i\geq1}$ be IID random variables such that
\[
\mathbb P(\omega_i=1)=\mathbb P(\omega_i=-1)=\frac12,
\qquad
\mathbb E[\omega_i]=0,
\qquad
\operatorname{Var}(\omega_i)=1.
\]

Define the scaled random walk:
\[
W^{(m)}(t)=\frac{1}{\sqrt m}\sum_{i=1}^{\lfloor mt\rfloor}\omega_i.
\]

\subsubsection*{Mean}

Using linearity of expectation:
\[
\mathbb E[W^{(m)}(t)]
=
\frac{1}{\sqrt m}\sum_{i=1}^{\lfloor mt\rfloor}\mathbb E[\omega_i]
=
\frac{1}{\sqrt m}\sum_{i=1}^{\lfloor mt\rfloor}0=0.
\]

\subsubsection*{Variance}

Since the $\omega_i$ are independent:
\[
\operatorname{Var}\!\left(W^{(m)}(t)\right)
=
\frac{1}{m}\sum_{i=1}^{\lfloor mt\rfloor}\operatorname{Var}(\omega_i)
=
\frac{1}{m}\lfloor mt\rfloor.
\]

As $m\to\infty$,
\[
\frac{\lfloor mt\rfloor}{m}\to t.
\]

Hence:
\[
\operatorname{Var}(W^{(m)}(t))\to t.
\]

\subsubsection*{Increment from $s$ to $t$}

For $0\le s<t$:
\[
W^{(m)}(t)-W^{(m)}(s)
=
\frac{1}{\sqrt m}\sum_{i=\lfloor ms\rfloor+1}^{\lfloor mt\rfloor}\omega_i.
\]

Therefore:
\[
\mathbb E[W^{(m)}(t)-W^{(m)}(s)]=0,
\]
and
\[
\operatorname{Var}(W^{(m)}(t)-W^{(m)}(s))
=
\frac{\lfloor mt\rfloor-\lfloor ms\rfloor}{m}\to t-s.
\]

\subsubsection*{Central Limit Theorem}

The increment is the normalized sum of many IID centered variables. By the Central Limit Theorem:
\[
W^{(m)}(t)-W^{(m)}(s)\xrightarrow[m\to\infty]{d}\mathcal N(0,t-s).
\]

\subsubsection*{Motivation for Brownian Motion}

This suggests defining a continuous-time limit process $\{W_t\}_{t\ge0}$ such that:

\begin{itemize}
    \item $W_0=0$,
    \item increments are Gaussian:
    \[
    W_t-W_s\sim\mathcal N(0,t-s),
    \]
    \item increments over disjoint intervals are independent,
    \item sample paths are continuous.
\end{itemize}

This process is called \textbf{Brownian motion}.

%---------------------------------------------------

\subsection*{2.2 Quadratic Variation}

Let
\[
0=t_0<t_1<\cdots<t_n=t
\]
be a partition of $[0,t]$, and define
\[
\Delta t_i=t_{i+1}-t_i.
\]

Set
\[
Q_n(t)=\sum_{i=0}^{n-1}(W_{t_{i+1}}-W_{t_i})^2.
\]

\subsubsection*{Expectation}

Since Brownian increments satisfy
\[
W_{t_{i+1}}-W_{t_i}\sim\mathcal N(0,\Delta t_i),
\]
we have:
\[
\mathbb E[(W_{t_{i+1}}-W_{t_i})^2]=\Delta t_i.
\]

Hence:
\[
\mathbb E[Q_n(t)]
=
\sum_{i=0}^{n-1}\Delta t_i=t.
\]

\subsubsection*{Representation with Standard Normal Variables}

Write
\[
W_{t_{i+1}}-W_{t_i}=\sqrt{\Delta t_i}\,Z_i,
\qquad Z_i\sim\mathcal N(0,1).
\]

Then:
\[
(W_{t_{i+1}}-W_{t_i})^2=\Delta t_i\, Z_i^2.
\]

Since $Z_i^2\sim\chi^2_1$, each squared increment is a time-step multiplied by a chi-squared random variable.

\subsubsection*{Variance of One Term}

Using $\operatorname{Var}(Z_i^2)=2$:
\[
\operatorname{Var}\!\left((W_{t_{i+1}}-W_{t_i})^2\right)
=
(\Delta t_i)^2\operatorname{Var}(Z_i^2)
=
2(\Delta t_i)^2.
\]

\subsubsection*{Variance of $Q_n(t)$}

Because Brownian increments are independent:
\[
\operatorname{Var}(Q_n(t))
=
\sum_{i=0}^{n-1}2(\Delta t_i)^2.
\]

Now:
\[
\sum_{i=0}^{n-1}(\Delta t_i)^2
\le
\left(\max_i \Delta t_i\right)\sum_{i=0}^{n-1}\Delta t_i
=
t\max_i\Delta t_i.
\]

Thus:
\[
\operatorname{Var}(Q_n(t))
\le
2t\max_i\Delta t_i \to 0
\quad\text{as}\quad \max_i\Delta t_i\to0.
\]

Since also $\mathbb E[Q_n(t)]=t$, we conclude:
\[
Q_n(t)\to t
\]
(in $L^2$, hence in probability).

\subsubsection*{Informal Differential Rule}

This justifies the It\^o rule:
\[
(dW_t)^2=dt.
\]

Moreover:

\[
dt^2=o(dt),\qquad dt\,dW_t=o(dt),
\]
because $dW_t$ is of size $\sqrt{dt}$, so:
\[
dt\,dW_t\sim dt^{3/2},\qquad dt^2\sim dt^2.
\]

Both are negligible compared with $dt$, while $(dW_t)^2\sim dt$ survives.

%---------------------------------------------------

\subsection*{2.3 Black--Scholes with Inventory Costs}

Assume:
\[
dS_t=\mu S_t\,dt+\sigma S_t\,dW_t.
\]

Let $V(S,t)$ be the derivative price.

\subsubsection*{It\^o's Lemma}

By It\^o:
\[
dV=V_t\,dt+V_S\,dS+\frac12V_{SS}(dS)^2.
\]

Since:
\[
(dS)^2=\sigma^2S^2dt,
\]
we get:
\[
dV=
\left(V_t+\mu SV_S+\frac12\sigma^2S^2V_{SS}\right)dt
+\sigma SV_S\,dW_t.
\]

\subsubsection*{Replicating Portfolio}

Let:
\[
X_t=\Delta_t S_t+\Phi_t B_t,
\qquad dB_t=rB_tdt.
\]

Equivalently:
\[
V_t-\Delta_tS_t=\Phi_tB_t.
\]

These are the same statement: the option value equals holdings in stock plus holdings in bond.

\subsubsection*{Inventory Cost}

Holding one unit of stock costs $cS_tdt$ per unit time, so:
\[
dX_t=\Delta_t\,dS_t+\Phi_t\,dB_t-c\Delta_tS_tdt.
\]

Substitute:
\[
dX_t=
\Delta_t(\mu S_tdt+\sigma S_tdW_t)+\Phi_trB_tdt-c\Delta_tS_tdt.
\]

Since:
\[
\Phi_tB_t=V-\Delta_tS,
\]
then:
\[
dX_t=
\Delta_t\mu Sdt+\Delta_t\sigma SdW_t+r(V-\Delta_tS)dt-c\Delta_tSdt.
\]

So:
\[
dX_t=
\left[rV+\Delta_tS(\mu-r-c)\right]dt+\Delta_t\sigma SdW_t.
\]

Now impose replication: $dX_t=dV$.

Matching diffusion terms:
\[
\Delta_t=V_S.
\]

Matching drift terms:
\[
V_t+\mu SV_S+\frac12\sigma^2S^2V_{SS}
=
rV+V_SS(\mu-r-c).
\]

Simplify (the $\mu$ terms cancel):
\[
V_t+(r+c)SV_S+\frac12\sigma^2S^2V_{SS}-rV=0.
\]

\[
\boxed{
V_t+(r+c)SV_S+\frac12\sigma^2S^2V_{SS}-rV=0
}
\]

\subsubsection*{Comparison with Dividend Yield}

Standard Black--Scholes with continuous dividend yield $q$:
\[
V_t+(r-q)SV_S+\frac12\sigma^2S^2V_{SS}-rV=0.
\]

Comparing coefficients:
\[
r-q=r+c
\quad\Longrightarrow\quad
q=-c.
\]

Thus inventory cost acts like a \textbf{negative dividend yield}.

%---------------------------------------------------

\subsection*{2.4 Black--Scholes when $\ln S_t$ is Ornstein--Uhlenbeck}

Assume:
\[
d(\ln S_t)=\kappa(\theta-\ln S_t)dt+\sigma dW_t.
\]

Let:
\[
X_t=\ln S_t,\qquad S_t=e^{X_t}.
\]

\subsubsection*{Dynamics of $S_t$}

Apply It\^o to $f(X)=e^X$:

\[
f'(X)=e^X=S,\qquad f''(X)=e^X=S.
\]

Hence:
\[
dS_t=f'(X_t)dX_t+\frac12f''(X_t)(dX_t)^2.
\]

Since $(dX_t)^2=\sigma^2dt$:
\[
dS_t
=
S_t\left[\kappa(\theta-\ln S_t)+\frac12\sigma^2\right]dt+\sigma S_t dW_t.
\]

\subsubsection*{Pricing PDE}

Write stock dynamics as:
\[
dS_t=\mu(S_t,t)S_tdt+\sigma S_tdW_t,
\]
where
\[
\mu(S,t)=\kappa(\theta-\ln S)+\frac12\sigma^2.
\]

Using the same replication argument as before, the drift cancels out. Therefore the PDE is:

\[
\boxed{
V_t+rSV_S+\frac12\sigma^2S^2V_{SS}-rV=0
}
\]

\subsubsection*{Comparison with Standard Black--Scholes}

This is the same PDE as the standard Black--Scholes equation (without dividends).

Although the \emph{physical} stock drift is mean-reverting in log-price, arbitrage pricing depends on hedging away risk, not on the real-world expected return. Once the stochastic term is replicated using the stock, the drift disappears and only $r$ remains.

%---------------------------------------------------

\subsection*{Conclusion}

The notation $o(dt)$ captures terms negligible relative to $dt$. In stochastic calculus, this matters because Brownian motion has nonzero quadratic variation:
\[
(dW_t)^2=dt.
\]
This is the key difference from ordinary calculus. It produces the second-derivative term in It\^o's lemma and therefore in pricing PDEs. Finally, even if the stock follows more complicated dynamics such as an Ornstein--Uhlenbeck log-price process, replication removes dependence on the physical drift, so the option pricing equation keeps the Black--Scholes structure.

\newpage
\section*{Exercise 4. Option Trading Strategies}

In this note we study six classical option strategies built from European calls, puts, stock, and cash.  
We assume the Black--Scholes framework and denote:

\[
\omega=T-t
\]

(time to maturity),

\[
C(S,t;K,T)
\]

the price of a European call with strike $K$ and maturity $T$, and

\[
P(S,t;K,T)
\]

the price of a European put with strike $K$ and maturity $T$.

For simplicity, we write:

\[
C(K):=C(S,t;K,T), \qquad P(K):=P(S,t;K,T).
\]

The key idea of the exercise is:

\[
\textbf{If a payoff is a sum of simpler payoffs, then its price is the same sum of prices.}
\]

This follows from no-arbitrage and linearity of replication.

%=========================================================
\subsection*{4.1 Risk-Neutral Pricing}

Under the risk-neutral probability measure $\mathbb Q$:

\[
V(S,t)=e^{-r(T-t)}\mathbb E^{\mathbb Q}[\Phi(S_T)\mid S_t=S].
\]

Therefore:

\[
V(S,t)=e^{-r\omega}\mathbb E^{\mathbb Q}[\Phi(S_T)\mid S_t=S].
\]

\subsubsection*{Why this works for any European payoff}

A European derivative pays only at time $T$, and its payoff depends only on $S_T$.  
Hence, once we know the distribution of $S_T$ under $\mathbb Q$, we can price any payoff $\Phi(S_T)$ by taking its discounted expectation.

\subsubsection*{Distribution of $S_T$}

Under Black--Scholes:

\[
dS_u=rS_u\,du+\sigma S_u\,dW_u.
\]

Applying It\^o's lemma to $\ln S_u$:

\[
d(\ln S_u)=\left(r-\frac12\sigma^2\right)du+\sigma dW_u.
\]

Integrating from $t$ to $T$:

\[
\ln S_T=\ln S+\left(r-\frac12\sigma^2\right)\omega+\sigma(W_T-W_t).
\]

Since:

\[
W_T-W_t\sim N(0,\omega),
\]

we obtain:

\[
\ln S_T\sim N\!\left(\ln S+\left(r-\frac12\sigma^2\right)\omega,\sigma^2\omega\right).
\]

Thus $S_T$ is lognormal.

\subsubsection*{Integral Form}

If $f_{LN}$ is the lognormal density:

\[
V(S,t)=e^{-r\omega}\int_0^\infty \Phi(x)\,f_{LN}(x)\,dx.
\]

%=========================================================
\subsection*{4.2 Strategy Decomposition}

We now analyze the six strategies listed in the assignment.

\bigskip
\textbf{Important notation.}

\begin{itemize}[leftmargin=1.8cm]
\item Long call = buy one call.
\item Short call = sell one call.
\item Long put = buy one put.
\item Long stock = buy one share.
\item Short two calls = sell two call options.
\end{itemize}

%=========================================================


\begin{center}
\small
\begin{longtable}{p{2.8cm} p{4.3cm} p{5.2cm} p{3.2cm}}
\toprule
\textbf{Strategy} & \textbf{Terminal Payoff $\Phi(S_T)$} & \textbf{How it is Built} & \textbf{Price at time $t$} \\
\midrule
\endfirsthead



Straddle
&
$(S_T-K)^+ + (K-S_T)^+$
&
Long 1 call with strike $K$ + Long 1 put with strike $K$
&
$C(K)+P(K)$
\\

Covered Call
&
$S_T-(S_T-K)^+$
&
Long stock + Short 1 call with strike $K$
&
$S-C(K)$
\\

Butterfly
&
$(S_T-K_1)^+-2(S_T-K_2)^++(S_T-K_3)^+$
&
Long call $K_1$, Short 2 calls $K_2$, Long call $K_3$
&
$C(K_1)-2C(K_2)+C(K_3)$
\\

Collar
&
$S_T+(K_P-S_T)^+-(S_T-K_C)^+$
&
Long stock + Long put $K_P$ + Short call $K_C$
&
$S+P(K_P)-C(K_C)$
\\

Protective Put
&
$S_T+(K-S_T)^+$
&
Long stock + Long put
&
$S+P(K)$
\\

Bull Call Spread
&
$(S_T-K_1)^+-(S_T-K_2)^+$
&
Long call $K_1$ + Short call $K_2$
&
$C(K_1)-C(K_2)$
\\

\bottomrule
\caption{Table 1: Summary of payoff, portfolio decomposition, and pricing formulas.}
\end{longtable}
\end{center}

%=========================================================
\subsection*{Explanation of the Main Constructions}

\subsubsection*{Straddle}

A straddle buys both a call and a put with the same strike.

\[
\Phi=(S_T-K)^+ + (K-S_T)^+
\]

If the stock goes far above $K$, the call pays.  
If the stock goes far below $K$, the put pays.  
Therefore the strategy benefits from \emph{large moves in either direction}.

Price:

\[
V=C(K)+P(K).
\]

%---------------------------------------------------------
\subsubsection*{Covered Call}

A covered call means owning the stock and selling a call.

\[
\Phi=S_T-(S_T-K)^+
\]

If $S_T\le K$, the call expires worthless and you keep the stock.  
If $S_T>K$, the upside is partially given away because the sold call is exercised.

Price:

\[
V=S-C(K).
\]

%---------------------------------------------------------
\subsubsection*{Butterfly}

The butterfly uses three strikes:

\[
K_1<K_2<K_3
\]

with $K_2$ usually centered.

Payoff:

\[
\Phi=(S_T-K_1)^+-2(S_T-K_2)^++(S_T-K_3)^+.
\]

It gains most when the final stock price stays near the middle strike $K_2$.

Price:

\[
V=C(K_1)-2C(K_2)+C(K_3).
\]

%---------------------------------------------------------
\subsubsection*{Collar}

A collar combines upside limitation and downside protection:

\[
\Phi=S_T+(K_P-S_T)^+-(S_T-K_C)^+.
\]

Interpretation:

\begin{itemize}
\item The put creates a floor (protection below $K_P$).
\item The short call caps gains above $K_C$.
\end{itemize}

Price:

\[
V=S+P(K_P)-C(K_C).
\]

%---------------------------------------------------------
\subsubsection*{Protective Put}

A protective put means buying stock and buying insurance.

\[
\Phi=S_T+(K-S_T)^+
\]

If the stock crashes, the put offsets losses below $K$.

Price:

\[
V=S+P(K).
\]

%---------------------------------------------------------
\subsubsection*{Bull Call Spread}

This strategy expresses a moderate bullish view.

\[
\Phi=(S_T-K_1)^-(S_T-K_2)^+
\]

(correct form:)

\[
\Phi=(S_T-K_1)^+-(S_T-K_2)^+, \qquad K_1<K_2.
\]

You gain when the stock rises, but profits are capped above $K_2$.

Price:

\[
V=C(K_1)-C(K_2).
\]

%=========================================================
\subsection*{4.3 Closed-Form Black--Scholes Prices}

We now compute each strategy explicitly.

The Black--Scholes formulas are:

\[
C(K)=SN(d_1)-Ke^{-r\omega}N(d_2),
\]

\[
P(K)=Ke^{-r\omega}N(-d_2)-SN(-d_1),
\]

where

\[
d_1(K)=\frac{\ln(S/K)+(r+\frac12\sigma^2)\omega}{\sigma\sqrt{\omega}},
\qquad
d_2(K)=d_1(K)-\sigma\sqrt{\omega}.
\]

When strikes differ, we compute $d_1,d_2$ using the corresponding strike.

%---------------------------------------------------------
\subsubsection*{Straddle}

By definition:

\[
V_{\text{str}}=C(K)+P(K).
\]

Insert the formulas:

\[
V_{\text{str}}
=
\big[SN(d_1)-Ke^{-r\omega}N(d_2)\big]
+
\big[Ke^{-r\omega}N(-d_2)-SN(-d_1)\big].
\]

Group stock terms and bond terms:

\[
V_{\text{str}}
=
S\big(N(d_1)-N(-d_1)\big)
+
Ke^{-r\omega}\big(N(-d_2)-N(d_2)\big).
\]

Using:

\[
N(-x)=1-N(x),
\]

we obtain:

\[
N(d_1)-N(-d_1)=2N(d_1)-1,
\]

\[
N(-d_2)-N(d_2)=1-2N(d_2).
\]

Therefore:

\[
\boxed{
V_{\text{str}}
=
S(2N(d_1)-1)+Ke^{-r\omega}(1-2N(d_2))
}
\]

%---------------------------------------------------------
\subsection*{Covered Call}

Price:

\[
V_{\text{cc}}=S-C(K).
\]

Substitute the call formula:

\[
V_{\text{cc}}
=
S-\big[SN(d_1)-Ke^{-r\omega}N(d_2)\big].
\]

Distribute the minus sign:

\[
V_{\text{cc}}
=
S-SN(d_1)+Ke^{-r\omega}N(d_2).
\]

Factor $S$:

\[
\boxed{
V_{\text{cc}}
=
S(1-N(d_1))+Ke^{-r\omega}N(d_2)
}
\]

%---------------------------------------------------------
\subsubsection*{Butterfly Spread}

Price:

\[
V_{\text{bf}}=C(K_1)-2C(K_2)+C(K_3).
\]

Now replace each call:

\[
V_{\text{bf}}
=
\big[S N(d_1^{(1)})-K_1e^{-r\omega}N(d_2^{(1)})\big]
\]

\[
-2\big[S N(d_1^{(2)})-K_2e^{-r\omega}N(d_2^{(2)})\big]
\]

\[
+\big[S N(d_1^{(3)})-K_3e^{-r\omega}N(d_2^{(3)})\big].
\]

Collect stock terms:

\[
S\Big(N(d_1^{(1)})-2N(d_1^{(2)})+N(d_1^{(3)})\Big),
\]

and discounted strike terms:

\[
-e^{-r\omega}
\Big(
K_1N(d_2^{(1)})-2K_2N(d_2^{(2)})+K_3N(d_2^{(3)})
\Big).
\]

Hence:

\[
\boxed{
V_{\text{bf}}
=
S\left[N(d_1^{(1)})-2N(d_1^{(2)})+N(d_1^{(3)})\right]
-
e^{-r\omega}
\left[
K_1N(d_2^{(1)})-2K_2N(d_2^{(2)})+K_3N(d_2^{(3)})
\right]
}
\]

where $d_i^{(j)}=d_i(K_j)$.

%---------------------------------------------------------
\subsubsection*{Collar}

Price:

\[
V_{\text{col}}=S+P(K_P)-C(K_C).
\]

Substitute:

\[
V_{\text{col}}
=
S+\big[K_Pe^{-r\omega}N(-d_2^P)-SN(-d_1^P)\big]
-\big[SN(d_1^C)-K_Ce^{-r\omega}N(d_2^C)\big].
\]

Expand:

\[
V_{\text{col}}
=
S-SN(-d_1^P)-SN(d_1^C)
+K_Pe^{-r\omega}N(-d_2^P)
+K_Ce^{-r\omega}N(d_2^C).
\]

Therefore:

\[
\boxed{
V_{\text{col}}
=
S\big[1-N(-d_1^P)-N(d_1^C)\big]
+
e^{-r\omega}
\big[
K_PN(-d_2^P)+K_CN(d_2^C)
\big]
}
\]

%---------------------------------------------------------
\subsubsection*{Protective Put}

Price:

\[
V_{\text{pp}}=S+P(K).
\]

Insert put formula:

\[
V_{\text{pp}}
=
S+Ke^{-r\omega}N(-d_2)-SN(-d_1).
\]

Factor $S$:

\[
\boxed{
V_{\text{pp}}
=
S\big(1-N(-d_1)\big)+Ke^{-r\omega}N(-d_2)
}
\]

Since $1-N(-d_1)=N(d_1)$, equivalently:

\[
\boxed{
V_{\text{pp}}=SN(d_1)+Ke^{-r\omega}N(-d_2)
}
\]

%---------------------------------------------------------
\subsubsection*{Bull Call Spread}

Price:

\[
V_{\text{bull}}=C(K_1)-C(K_2).
\]

Substitute:

\[
V_{\text{bull}}
=
\big[S N(d_1^{(1)})-K_1e^{-r\omega}N(d_2^{(1)})\big]
-
\big[S N(d_1^{(2)})-K_2e^{-r\omega}N(d_2^{(2)})\big].
\]

Expand:

\[
V_{\text{bull}}
=
S\big(N(d_1^{(1)})-N(d_1^{(2)})\big)
-
e^{-r\omega}
\big(K_1N(d_2^{(1)})-K_2N(d_2^{(2)})\big).
\]

Thus:

\[
\boxed{
V_{\text{bull}}
=
S\left[N(d_1^{(1)})-N(d_1^{(2)})\right]
-
e^{-r\omega}
\left[K_1N(d_2^{(1)})-K_2N(d_2^{(2)})\right]
}
\]

%=========================================================
\subsection*{4.4 Greeks}

Because derivatives are linear operators:

\[
G(aV_1+bV_2)=aG(V_1)+bG(V_2),
\]

for any Greek $G$.

So the Greeks of each strategy are obtained by summing the Greeks of its components.

%=========================================================


\begin{center}
\small
\begin{longtable}{p{3cm} p{11cm}}
\toprule
\textbf{Strategy} & \textbf{Greeks} \\
\midrule
\endfirsthead

\toprule
\textbf{Strategy} & \textbf{Greeks} \\
\midrule
\endhead

Straddle &
$\Delta=\Delta_C+\Delta_P,\quad
\Gamma=\Gamma_C+\Gamma_P,\quad
\Theta=\Theta_C+\Theta_P,\quad
\nu=\nu_C+\nu_P,\quad
\rho=\rho_C+\rho_P$
\\

Covered Call &
$\Delta=1-\Delta_C,\quad
\Gamma=-\Gamma_C,\quad
\Theta=-\Theta_C,\quad
\nu=-\nu_C,\quad
\rho=-\rho_C$
\\

Butterfly &
$\Delta=\Delta_1-2\Delta_2+\Delta_3$

$\Gamma=\Gamma_1-2\Gamma_2+\Gamma_3$

$\Theta=\Theta_1-2\Theta_2+\Theta_3$

$\nu=\nu_1-2\nu_2+\nu_3$

$\rho=\rho_1-2\rho_2+\rho_3$
\\

Collar &
$\Delta=1+\Delta_P-\Delta_C$

$\Gamma=\Gamma_P-\Gamma_C$

$\Theta=\Theta_P-\Theta_C$

$\nu=\nu_P-\nu_C$

$\rho=\rho_P-\rho_C$
\\

Protective Put &
$\Delta=1+\Delta_P,\quad
\Gamma=\Gamma_P,\quad
\Theta=\Theta_P,\quad
\nu=\nu_P,\quad
\rho=\rho_P$
\\

Bull Call Spread &
$\Delta=\Delta_1-\Delta_2$

$\Gamma=\Gamma_1-\Gamma_2$

$\Theta=\Theta_1-\Theta_2$

$\nu=\nu_1-\nu_2$

$\rho=\rho_1-\rho_2$
\\

\bottomrule
\caption{Greeks obtained by linearity.}
\end{longtable}
\end{center}

%=========================================================
\newpage
\subsection*{4.5 Interpretation}

\begin{center}
\small
\begin{longtable}{p{2.8cm} p{3.3cm} p{3.2cm} p{5.2cm}}
\toprule
\textbf{Strategy} & \textbf{Market View} & \textbf{Gamma Profile} & \textbf{Typical Use Case} \\
\midrule
\endfirsthead

\toprule
\textbf{Strategy} & \textbf{Market View} & \textbf{Gamma Profile} & \textbf{Typical Use Case} \\
\midrule
\endhead

Straddle &
Volatility-based &
Long gamma &
Expecting a strong move but unsure of direction
\\

Covered Call &
Mild bullish / neutral &
Short gamma &
Generate extra income on an existing stock position
\\

Butterfly &
Range-bound &
Mixed &
Expecting price near a target level
\\

Collar &
Protective bullish &
Mixed / limited &
Keep stock exposure while reducing downside risk
\\

Protective Put &
Protective bullish &
Long gamma &
Insurance for a long-term stock investor
\\

Bull Call Spread &
Moderately bullish &
Positive but capped &
Cheaper bullish position than buying a naked call
\\

\bottomrule
\caption{Interpretation of the six strategies.}
\end{longtable}
\end{center}

%=========================================================
\subsection*{Conclusion}

The whole exercise is based on one fundamental principle:

\[
\boxed{\text{Complex payoffs can be decomposed into simple payoffs.}}
\]

Once we know how to price calls, puts, stock, and bonds under Black--Scholes, every strategy follows immediately by linearity. The same reasoning applies to sensitivities (Greeks): portfolio risk is simply the sum of the risks of each component.

\end{document}
